% 主题切换示例：同一份正文，只改 \usetheme 一行即可分别编译。
% 变体入口见 tests/switch-hitminimalist.tex、tests/switch-hit.tex、tests/switch-classic.tex、tests/switch-madrid.tex；
% 说明：hit 为默认 touying 主题，classic 为经典主题，minimalist 为极简主题，Madrid 为内建主题。
% 它们只定义 \hitdemo，正文部分完全一致。
\documentclass[aspectratio=169]{ctexbeamer}
\providecommand{\hitdemo}{minimalist}
\ifnum\pdfstrcmp{\hitdemo}{Madrid}=0
  \usetheme{Madrid}
\else\ifnum\pdfstrcmp{\hitdemo}{classic}=0
  \usetheme[classic]{hit}
\else\ifnum\pdfstrcmp{\hitdemo}{touying}=0
  \usetheme{hit}
\else\ifnum\pdfstrcmp{\hitdemo}{hit}=0
  \usetheme{hit}
\else\ifnum\pdfstrcmp{\hitdemo}{minimalist}=0
  \usetheme[minimalist]{hit}
\else\ifnum\pdfstrcmp{\hitdemo}{hitminimalist}=0
  \usetheme[minimalist]{hit}
\else\ifnum\pdfstrcmp{\hitdemo}{hitacademic}=0
  \usetheme[minimalist]{hit}
\else
  \usetheme[\hitdemo]{hit}
\fi\fi\fi\fi\fi\fi\fi
\usepackage{beamercmdhit}

\title[同一份正文]{同一份正文下的主题切换}
\englishtitle{One Body, Three Themes}
\titlecontext{公共接口一致性示例}
\author{汇报人：待填写}
\institute{哈尔滨工业大学}
\date{汇报日期：待填写}

\begin{document}

\TitleSlide
\OutlineSlide

\section{公共接口}

\begin{frame}{各主题都支持的内容}
  \subhead{正文与小标题}
  这份正文只使用 \texttt{\textbackslash usetheme} 切换主题，其余各行不变。

  \insight{\texttt{\textbackslash insight} 使用主题配色，不依赖某个主题的私有颜色。}
  \notebox{\texttt{\textbackslash notebox} 在各主题下都保留 2 pt 左侧竖线。}
  \captiontext{\texttt{\textbackslash captiontext} 使用主题的淡色文字。}
\end{frame}

\begin{frame}{指标与条目}
  \begin{columns}[T,onlytextwidth]
    \begin{column}{.30\textwidth}
      \metric{\linewidth}{1.8×}{解析模型给出的加速比}
    \end{column}
    \begin{column}{.64\textwidth}
      \begin{denseitems}
        \item 同一条公共命令只定义一次。
        \item 视觉参数由主题的版式槽位给出。
        \item 内建主题下退回通用实现。
      \end{denseitems}
    \end{column}
  \end{columns}
\end{frame}

\section{整页命令}

\begin{frame}{整页命令的兼容调用}
  \begin{itemize}
    \item \texttt{\textbackslash TitleSlide}：各主题各自的封面。
    \item \texttt{\textbackslash OutlineSlide}：自动读取本节之前的正文主章节。
    \item \texttt{\textbackslash FocusSlide\{内容\}}：单参数调用保持可用。
    \item \texttt{\textbackslash EndSlide\{内容\}}：结束页。
  \end{itemize}
\end{frame}

\FocusSlide{公共接口只维护一份}

\EndSlide{感谢聆听\quad 欢迎讨论}

\end{document}
